\documentclass[output=paper]{langscibook}

\author{
    Sabine Laszakovits\affiliation{Austrian Academy of Sciences; University of Connecticut} and
    {Ronice Müller de} Quadros\affiliation{Universidade Federal de Santa Catarina} and 
    {Emily Jo} Noschese\affiliation{University of Hawai'i} and 
    Diane Lillo-Martin\affiliation{University of Connecticut; Haskins Laboratories}
}
\title{Object shift in ASL and Libras}
\abstract{ASL and Libras have an object-shift construction by which the canonical SVO order is changed to SOV. In both sign languages, this ordering is mandatory for V marked with durative/continuative aspect (reduplicated movement), optional for V that agrees with O in locus, and not allowed with plain V. When V agrees with O in handshape, ASL requires OV ordering whereas Libras allows both OV and VO ordering. We present an analysis that derives these data with a combination of syntactic movement of O and violable, equally-ranked PF-constraints as proposed by \citet{BW.2012}. Unlike \citeapos{Matsuoka.1997} and \citeapos{Braze.2004} proposals, we do not move V to a head on the right in violation of the Final-over-final constraint    \citep{BiberauerHR.2014,SheehanBRH.2017}.}

\begin{document}
\lehead{Laszakovits et al.}
\maketitle

% ===================================================================
\section{Introduction}
\label{lasz:sec:1}

The underlying word order in modern American Sign Language (ASL)
and Brazilian sign language (Língua Brasileira de Sinais, Libras) has
often been argued to be subject-verb-object (SVO; 
see \citealp{Fischer.1975,Liddell.1980,Padden.1988} for ASL; 
\citealp{Quadros.1999,Quadros.2003} for Libras), 
which we adopt here. However, in both languages, word order
variations are possible. In this paper, we focus on the construction with
the word order SOV, which has been termed ``object shift''. This
construction differs from topicalization structures like O'SV, VO'S,
S'VO, or S'O'V in that the latter contain a prosodic break after the
topic (indicated here by an apostrophe) and may contain nonmanual
marking (raised brows and a slight upward chin tilt; 
\citealp{Liddell.1977,Liddell.1980,Padden.1988}) during the utterance of the topic.%
\footnote{\citet{Fischer.1990} argues that object-shift constructions in ASL involve
    ``mini-topicalization'' of O to some intermediate specifier position. She corroborates
    this analysis by reporting a definiteness effect found in SOV orders, which is typical
    of topics. However, unlike typical topicalizations, SOV constructions lack a prosodic
    break after O and cannot mark O with nonmanual topicalization marking. Our
    analysis is similar to \citeauthor{Fischer.1990}'s in that we think O moves to some intermediate
    specifier position. However, we have not been able to reproduce this definiteness
    effect for ASL, and we see no reason to call this movement topicalization.
}
We restrict our attention to verbs that take a direct object (such as `buy X') and
possibly an indirect object (such as `send X to Y'), but do not discuss
verbs that denote movement to or from a location (such as `put X on Z').%
\footnote{Furthermore, we distinguish object-shift from locative constructions of the
    form location-subject-predicate, which have been taken as instances of the word
    order OSV because, like object-shift, they lack the prosodic break of topicalization
    structures \citep{Liddell.1980}. Locative constructions differ from object-shift
    constructions at least in the following properties: 
    \begin{inparaenum}[(i)]
        \item the word order is OSV (vs.~SOV),
        \item the predicate is restricted to `is located at' (vs.~any verb), and 
        \item there is no minimally different derivation with the word order SVO (which there is for
    object-shift; see e.g.~\REF{lasz:ex:4}--\REF{lasz:ex:7} and throughout the paper).
    \end{inparaenum}
}

There are three main triggers for object-shift in ASL and
Libras: 
\begin{inparaenum}[1.]
    \item durative/continuative aspect on the verb, which makes object-shift obligatory; 
    \item verbal agreement in handshape with the object, which makes object-shift 
    optional in Libras and obligatory in ASL; and 
    \item verbal agreement in locus with the object, which makes object-shift optional in both languages. 
\end{inparaenum}
The effect of durative aspect on word order is illustrated in \REF{lasz:ex:1}--\REF{lasz:ex:2}.%
\footnote{
    Following standard practice in sign linguistics, signs are glossed using
    English words in all caps. Most of our examples represent sign sequences that are
    grammatical in both ASL and Libras, though the actual signs are of course different.
    `IX' is the gloss used for a pointing indexical sign which serves
    pronominal functions. `\laszAsp{}' indicates an aspectual marker involving reduplication
    of the sign root. `\laszHs{\_}' indicates the use of a classifier handshape. `\laszLoc{\_}'
    indicates that a sign is produced using a spatial locus other than default (`a', `b',
    etc. indicate distinct loci but not physical location). `\laszDir{\_}{\_}' indicates that a sign
    moves from one locus to another. \par 
    In the English translation, past tense is used although neither ASL nor
    Libras has grammatical tense marking; similarly, we often use `a/the' for nouns
    because in both of these sign languages either translation is possible.
}%
\footnote{When a source is not provided, the judgments primarily come from the
    2\textsuperscript{nd} and 3\textsuperscript{rd} co-authors, who are native signers of Libras and ASL respectively.
}
\ea
    \label{lasz:ex:1}
    \ea[]{
        \label{lasz:ex:1a}
        \LaConcept{IX1 WINE DRINK}\laszAsp{}. 
        \hfill 
        \cmark ASL, \cmark Libras \\ 
        `I drank wine continuously.'
    }
    \ex[*]{
        \label{lasz:ex:1b}
        IX1 DRINK\laszAsp{} WINE. 
        \hfill 
        *ASL, *Libras
    }
    \z 
\ex
    \label{lasz:ex:2}
    \ea[]{
        \label{lasz:ex:2a}
        \LaConcept{MY} \LaConcept{SISTER} \LaConcept{LETTER} \LaConcept{SEND}\laszAsp{}. 
        \hfill 
        \cmark ASL, \cmark Libras \\ 
        `My sister repeatedly sent a/the letter(s).'
    }
    \ex[*]{ 
        \label{lasz:ex:2b}
        \LaConcept{MY} \LaConcept{SISTER} \LaConcept{SEND}\laszAsp{} \LaConcept{LETTER}. 
        \hfill 
        *ASL, *Libras 
    }
    \z 
\z
Verbal agreement can target the object's noun-class%
\footnote{%
    In this sense we follow \citeapos{Benedicto.Brentari.2004} proposal, whereby
    classifiers are treated as a type of agreement between V and O/S.
}, 
whereby V
changes its handshape (`hs') to a classifier as illustrated in~\REF{lasz:ex:3}, or it
can target the object's referent, whereby V changes the end-point of its
movement (`dir') to the object's locus (`loc'), as seen in~\REF{lasz:ex:4}.%
\footnote{
    \citet{Fischer.1975} proposed that object-shift is also licensed if the relation
    between V and O is established in the semantics, such that SOV is possible only if it
    is clear from the semantics of V, S, and O which argument is the agent and which is
    the patient. However, \citet{Liddell.1980} showed that this approach makes incorrect
    predictions, and that syntactic agreement is a better way to go.
}
\ea 
    \label{lasz:ex:3}
    The handshape of `\textsc{give}' changes to the \laszClaw{} handshape indicative of handling an apple. 
    \ea[]{ 
        \LaConcept{SALLY} APPLE \textsc{give}\laszHs{\laszClaw}. 
        \hfill 
        \cmark ASL, \cmark Libras 
    }
    \ex[]{ 
        \LaConcept{SALLY} \textsc{give}\laszHs{\laszClaw} APPLE. 
        \hfill
        \cmark Libras 
    }
    \ex[*]{
        \LaConcept{SALLY} \textsc{give}\laszHs{\laszClaw} APPLE. 
        \hfill 
        *ASL \\ 
        all: `Sally gave someone a/the apple.'
    }
    \z 
\ex
    \label{lasz:ex:4}
    The direction of `HELP' changes to the locus of `ANA'. 
    \ea 
        IX\laszLoc{a} MARIA\laszLoc{a} IX\laszLoc{b} ANA\laszLoc{b} HELP\laszDir{a}{b}. \\
        \citep[from][6]{Quadros.etal.2004}
        \hfill 
        \cmark Libras, \cmark ASL 
    \ex 
        IX\laszLoc{a} MARIA\laszLoc{a} HELP\laszDir{a}{b} IX\laszLoc{b} ANA\laszLoc{b}. \\
        \citep[from][5]{Quadros.etal.2004} 
        \hfill 
        \cmark Libras, \cmark ASL \\
        both: `Maria helped Ana.'
    \z 
\z 
For some verbs, locus agreement involves signing the verb at the
object's locus as illustrated in \REF{lasz:ex:5} (termed ``locationality'' by \citealp{Fischer.Gough.1978}, 
``spatialization'' by \citealp{Quadros.etal.2004}, 
and ``co-localization'' by \citealp{Lourenco.Wilbur.2018}; 
see also \citealp{Bergman.1980,Liddell.1980,Costello.2015,Smith.1990}, i.a.; ASL \& Libras),
or by adding an auxiliary as in \REF{lasz:ex:6} (Libras only; \citealp{Quadros.1999}). We
will treat all of these options as locus-agreement strategies and
postulate a uniform syntax.
\ea 
    \label{lasz:ex:5}
    Spatialized plain verbs
    \ea 
        MAN BICYCLE\laszLoc{a} \LaConcept{BUY}\laszLoc{a}. 
        \hfill 
        \cmark ASL, \cmark Libras \\ 
        \citep[from][9]{Quadros.etal.2004}
    \ex 
        MAN \LaConcept{BUY}\laszLoc{a} BICYCLE\laszLoc{a}. 
        \hfill 
        \cmark ASL, \cmark Libras \\ 
        \citep[from][9]{Quadros.etal.2004} \\ 
        both: `The man bought a/the bicycle.'
    \z 
\ex 
    \label{lasz:ex:6}
    Auxiliary 
    \ea 
        IX\laszLoc{a} JOÃO\laszLoc{a} IX\laszLoc{b} MARIA\laszLoc{b} AUX\laszDir{a}{b} SUPPORT. 
        \hfill 
        \cmark Libras 
        \\ 
        `Joao supported Maria.'
        \citep[from][7]{Quadros.etal.2004} 
    \ex 
        IX\laszLoc{a} MARIA\laszLoc{a} SUPPORT IX\laszLoc{b} LULA\laszLoc{b}. 
        \hfill 
        \cmark Libras 
        \\ 
        `Maria supported Lula.' 
        \citep[from][5]{Quadros.etal.2004}
    \z 
\z 
The object-shift construction is not possible if neither trigger is
present, i.e., when the verb is ``plain'', as in~\REF{lasz:ex:7}.
\ea 
    \label{lasz:ex:7}
    \ea[*]{
        MAN NUMBER FORGET\laszPlain{}. 
        \hfill 
        *ASL
        % \\ 
        % \citep[from][]{Liddell.1980}
    }
    \ex[]{
        MAN FORGET\laszPlain{} NUMBER. 
        \hfill 
        \cmark ASL 
        \\ 
        `The man forgot a/the number.'
        \citep[from][]{Liddell.1980}
    }
    \ex[*]{
        IX JOHN SOCCER LIKE\laszPlain{}. 
        \hfill 
        *Libras
        % \\ 
        % \citep[from][61]{Quadros.1999}
    }
    \ex[]{
        IX JOHN LIKE\laszPlain{} SOCCER. 
        \hfill 
        \cmark Libras 
        \\ 
        `John likes soccer.'
        \citep[from][61]{Quadros.1999}
    }
    \z 
\z 
Table~\ref{lasz:tab:1} presents a summary of the acceptability judgments for the
basic order SVO and for the object-shift construction SOV for both
ASL and Libras.

\begin{table}
    \centering 
    \begin{tabular}{r c c c c}
        \lsptoprule 
        & V\laszPlain{} & V\laszLoc{\_} & V\laszHs{\_} & V\laszAsp{} \\\midrule
        ASL & 
            \laszTabColA{} & 
            \laszTabColB{} & 
            \laszTabColC{} & 
            \laszTabColC{} 
            \\ 
        Libras & 
            \multirow{-2}{*}{\laszTabColA{} \shortstack{\cmark SVO \\ \xmark SOV}} & 
            \multirow{-2}{*}{\laszTabColB{} \shortstack{\cmark SVO \\ \cmark SOV}} & 
            \laszTabColB{} & 
            \multirow{-2}{*}{\laszTabColC{} \shortstack{\xmark SVO \\ \cmark SOV}} 
            \\
        \lspbottomrule
    \end{tabular}
    \caption{Comparison of judgments for SVO and OSV}
    \label{lasz:tab:1}
\end{table}

Several proposals have been made regarding the syntax of object-shift
in ASL resp.~Libras. The difficulty in providing a satisfactory analysis of
this phenomenon seems to stem from the disparate nature and the
disparate obligatoriness of the constructions that trigger object shift.

For durative-aspect marking of V, \citet{Matsuoka.1997} and \citet{Braze.2004} 
have proposed that object shift arises because aspectual
morphology sits on an Asp head, and V must move there, and AspP is
right-headed. In section~\ref{lasz:sec:3} we will show that there is an empirical and a
conceptual problem with this analysis.

An alternative proposal is based on the fact that durative-aspect
marking is expressed by reduplication of the verb’s movement 
\citep{Klima.Bellugi.1979}. The complexity of a sign’s movement corresponds to
its phonological heaviness \citep{Brentari.1998}, and we know from several
spoken languages that there is a preference for phonologically heavy
elements to appear late in a sentence \parentext{\citealp{Ross.1967,Williams.2003}, i.a.}.
\citet{Brentari.1998} has argued that this preference also applies in the
object-shift construction with heavy verbs, but remains agnostic
whether the inversion is realized by V moving to the right, or by O
moving to the left triggered by V’s heaviness, as \citet{Liddell.1980}
proposed.%
\footnote{
    Note that according to \citeapos{Brentari.1998} analysis, locus-agreement per se
    does not constitute `complex movement', neither if it modifies the locus in which V
    is signed, nor if it modifies the locus of V’s endpoint or starting point. In section 6.5
    she discusses how phonological weight units interact with the order of syntactic
    constituents. In \citeauthor{Brentari.1998}’s model, movement corresponds to prosodic features, and
    each branching prosodic feature constitutes a weight unit. In ASL, \textsc{give} usually
    contains only a “direction”-movement, that is one weight unit. Both VO and OV
    ordering are allowed. When habitual aspect is added, \textsc{give} contains an additional
    “trill” movement, a total of two weight units. OV ordering is preferred over VO.
    Exhaustive aspect adds both a “trill” movement and an “arc” movement, totallying
    three weight units. Now, due to V’s heaviness, VO is unacceptable.
}

However, as we have seen, object shift is also related to
agreement between V and O, both noun-class (classifier) agreement
\REF{lasz:ex:3}; and referential (locus) agreement \REF{lasz:ex:4}--\REF{lasz:ex:5}. 
These agreement-data
are difficult to unify with the aspect-data. If one follows a syntactic
proposal for the aspect-data like \citet{Matsuoka.1997} or \citet{Braze.2004},
one has to explain why aspect-marking makes the OV order
obligatory, while agreement-marking allows both the OV and the VO
order. On the other hand, if one follows a phonology-driven proposal
for the aspect-data like \citet{Liddell.1980} or \citet{Brentari.1998}, one needs an
independent account for the agreement-data, because changes to V’s
direction or locus or handshape do not influence V’s heaviness
\citep{Brentari.1998}.

Should the aspect-construction receive an analysis independent
of the agree\-ment-con\-struc\-tion, then, or is there a way to unify them?
We will propose a unified analysis that takes into account both
phonological heaviness and syntactic agreement projections. We use
\citeapos{Bobaljik.Wurmbrand.2012} proposal that, for any given
derivation and its resulting LF, a combination of inviolable and
violable, but equally ranked, PF-constraints determine the best
realization of this LF. We will introduce four independently attested
PF-constraints and show that they derive the data presented here
perfectly. In this way, we demonstrate that object shift in ASL and
Libras arises from the interaction of syntax and PF, and we strengthen
current assumptions about the size of the verbal domain by making
explicit the postulated projections.

The rest of the chapter is structured as follows: In section~\ref{lasz:sec:2}, we
present our analysis in detail and show that it makes the correct
predictions. In section~\ref{lasz:sec:3}, we discuss the alternative proposals by
\citet{Fischer.Janis.1992}, \citet{Matsuoka.1997}, and \citet{Braze.2004}, who
argue for rightward-movement of V. Section~\ref{lasz:sec:4} concludes.

% ===================================================================
\section{Analysis}
\label{lasz:sec:2}
    
\subsection{Phrase structure}
\label{lasz:sec:21}

Our analysis extends the proposals made by \citet{Quadros.etal.2004}, 
\citet{Quadros.DLM.2010}, and \citet{Gokgoz.2013}, who follow the
common analysis of object shift in spoken languages (\citealp{Holmberg:1986}
and much work since) and suggest that in object-shift constructions,
the order of V and O is inverted by movement of O to the left. Further,
we adopt the following insights from the literature:

As stated before, we follow \citet{Fischer.1975}, \citet{Liddell.1980},
\citet{Quadros.1999} and many others in assuming that the underlying order
of ASL and Libras is SVO. In other words, VP is left-headed.

Regarding agreement with noun-classes, we adopt \citeapos{Benedicto.Brentari.2004}
proposal: For classifier-agreement with O, there is
a dedicated phrase ClassOP immediately dominating VP, where the
handshape-feature on V is set by V moving to the head of ClassOP and O
moving to Spec,ClassOP, followed by Agree between O and V.%
\footnote{
    The feature value that O shares with V is not the handshape feature of O’s
    sign, but O’s sign’s noun-class resp.~information about the shape of O’s referent.
}
If the
classifier expressed on V is a handling classifier, there is a second,
higher ClassSP responsible for agreement between V and S \citep{Benedicto.Brentari.2004}.

For agreement with loci, we adopt the proposal that loci are the
realization of morphosyntactic features, similar to the proposals by
\citet{Neidle.etal.2000}, \citet{Kuhn.2016}, and \citet{Pfau.etal.2018}. We call the
phrase where V agrees with O in locus AgrOP, and assume that it
dominates ClassOP or, in its absence, VP.%
\footnote{
    \citet{Gokgoz.2013} argues that the phrase that objects shift to is in the
    TP-domain. He elicited word order data where in SOV sentences in ASL, manual
    sentential negation in preverbal position is rated better between O and V (SONegV)
    than before O (SNegOV): Assuming that manual negation appears in a NegP and
    that this NegP marks the edge of the VP-domain, O must raise beyond.
    
    We note that in his Experiment 3, the most accepted placement of manual
    negation is in sentence-final position for plain V, classifier-agreeing V,
    locus-agreeing V, and ditransitives agreeing in classifier and locus. (He does not
    discuss durative/continuative aspect.) We tentatively suggest that the realization of
    manual (and non-manual) negation may follow similar PF-constraints as the
    realization of O as discussed in this paper, but we consider an investigation of
    negation outside the scope of the present proposal.
    
    \citeauthor{Gokgoz.2013} also presents examples where shifted O has a presuppositional
    reading, which indicates, in line with \citet{Diesing1992}, that O moves out of the VP.
    This presuppositional effect is also observed in \citet{Napoli.SS.RMQ.2017}. We have not
    collected further data on presuppositional readings of O, but we did test for a
    definiteness effect of O in ASL (parallel to Scandinavian object-shift; \citealp{Holmberg:1986})
    and did not find one.
}
Similarly, for
directional verbs that contain movement from one locus to another, the
endpoint referring to S is established by agreement between V and S
in a projection AgrSP.

For aspect-marked verbs, we adopt parts of \citeapos{Matsuoka.1997}
and \citeapos{Braze.2004} analysis of the syntax of aspect: When the verb
displays durative-aspect agreement (reduplication), V moves to an Asp
head. However, we differ from these proposals in the headedness of
AspP: We contend that Asp is left-headed, like all verbal projections in
ASL and Libras. This allows our proposal to be consistent with the
Final-over-Final Constraint (FOFC; \citealp{BiberauerHR.2014,SheehanBRH.2017}) 
as will be discussed in section~\ref{lasz:sec:33}.

The tree in Fig.~\ref{lasz:fig:1} illustrates our assumptions about the phrase
structure of ASL and Libras for transitive verbs. Regarding
ditransitives like \textsc{give}, which agree with the direct object in
handshape, but with the indirect object in locus, we tentatively assume
that the phrase that introduces the indirect object is projected between 
ClassOP and AgrOP. In this way, ClassO will probe down and find O,
while AgrO will find IO because it is structurally higher than O.%
\footnote{
    We follow \citet{Quadros.etal.2004} in the assumption that the projections
    ClassOP, AgrOP, AspP, and ClassSP are only present when the verb shows the respective marking,
    but nothing in our analysis hinges on this.
}%
\footnote{
    See also the discussion in \citet{Pfau.etal.2018} of \citeapos{Nevins.2011} proposal.
}

\begin{figure}
    \begin{tikzpicture}
        \Tree
        [.TP
            \node (SpecTP) {$\langle$S$\rangle$};
            [
                [.T ]
                [.ClassSP
                    \node (SpecClSP) {$\langle$S$\rangle$};
                    [
                        [.ClassS ]
                        [.AspP
                            \node (SpecAspP) {$\langle$O$\rangle$}; 
                            [
                                [.Asp ]
                                [.\textit{v}P
                                    \node (Sbase) {$\langle$S$\rangle$};
                                    [
                                        [.\textit{v} ]
                                        [.AgrOP
                                            \node (SpecAgrOP) {$\langle$O$\rangle$};
                                            [
                                                [.AgrO ]
                                                [.ClassOP
                                                    \node (SpecClOP) {$\langle$O$\rangle$};
                                                    [
                                                        [.ClassO ]
                                                        [.VP
                                                            [.V ]
                                                            \node (Obase) {$\langle$O$\rangle$};
                                                        ]
                                                    ]
                                                ]
                                            ]
                                        ]
                                    ]
                                ]
                            ]
                        ]
                    ]
                ]
            ]
        ]
        \draw[->] (Obase) [in=-90,out=-135]  to (SpecClOP);
        \draw[->] (SpecClOP) [in=-90,out=-135]  to (SpecAgrOP);
        \draw[->] (SpecAgrOP) [in=-90,out=-135]  to (SpecAspP);
        \draw[->] (Sbase) [in=-90,out=-135]  to (SpecClSP);
        \draw[->] (SpecClSP) [in=-90,out=-135]  to (SpecTP);
    \end{tikzpicture}
    \caption{Phrase structure of ASL and Libras}
    \label{lasz:fig:1}
\end{figure}

\subsection{Verbal agreement}
\label{lasz:sec:22}

We make the following assumptions about verbal agreement:
In sentences where V agrees in handshape/locus/aspect, V
moves to the respective syntactic head (Class/Agr/Asp). This movement
is covert. In other words, there is an inviolable constraint at PF stating
that it is always the lowest copy of V that will be realized. This is
similar to English, where V agrees with T, but fully lexical verbs are
never realized in T. Another example of such an inviolable constraint
is in \citeapos{Bobaljik.Wurmbrand.2012} analysis of Quantifier Raising in
English: English syntax has scrambling just like German and
Japanese, and this movement is visible at LF in all languages, but in
contrast to German and Japanese, English PF never pronounces the
higher copies, giving rise to QR.

That V does not move overtly is corroborated by the fact that
we never see V move past manual negation~\REF{lasz:ex:8} or adverbs of
frequency~\REF{lasz:ex:9} in ASL and Libras:

\ea
    \label{lasz:ex:8}
    \ea[]{ 
        \LaConcept{IX}\laszLoc{a} \LaConcept{MARIA}\laszLoc{a} \LaConcept{\textbf{NOT}} \LaConcept{HELP}\laszDir{a}{b} \LaConcept{IX}\laszLoc{b} \LaConcept{JOHN}\laszLoc{b}. 
        \\
        `Maria did not help John.'
        \hfill         
        \cmark Libras, 
        \cmark ASL 
        \\
        \citep[][\pnfmt{6}]{Quadros.etal.2004}
    }
    \ex[*]{ 
        \LaConcept{IX}\laszLoc{a} \LaConcept{MARIA}\laszLoc{a} \LaConcept{HELP}\laszDir{a}{b} \LaConcept{\textbf{NOT}} \LaConcept{IX}\laszLoc{b} \LaConcept{JOHN}\laszLoc{b}.
        \\
        \citep[][\pnfmt{6}]{Quadros.etal.2004}
        \hfill 
        *Libras, 
        *ASL 
    }
    \z 
\ex 
    \label{lasz:ex:9}
    \ea[]{
        \LaConcept{IX}\laszLoc{a} \LaConcept{MARIA}\laszLoc{a} \textbf{\LaConcept{SOMETIMES}} \LaConcept{HELP}\laszDir{a}{b} \LaConcept{IX}\laszLoc{b} \LaConcept{JOHN}\laszLoc{b}. 
        \hfill 
        \cmark Libras, 
        \cmark ASL
        \\
        `Maria sometimes helped John.'
        \citep[][\pnfmt{6}]{Quadros.etal.2004}
    }
    \ex[*]{
        \LaConcept{IX}\laszLoc{a} \LaConcept{MARIA}\laszLoc{a} \LaConcept{HELP}\laszDir{a}{b} \textbf{\LaConcept{SOMETIMES}}  \LaConcept{IX}\laszLoc{b} \LaConcept{JOHN}\laszLoc{b}. 
        \hfill 
        \cmark Libras, 
        \cmark ASL
        \\
        \citep[][\pnfmt{6}]{Quadros.etal.2004}
    }
    \z
\z 

\noindent 
When V moves to agree, this triggers syntactic movement of O to the
respective specifier position. That is, when V agrees in handshape
with O, O moves to the specifier of ClassOP. When V agrees in
locus with O, O moves to the specifier of AgrOP. When V agrees in
aspect, O moves to the specifier of AspP. Whether this movement of O
is overt or covert, i.e. which copy of O gets realized, depends on
PF-constraints as we will discuss below. When agreement doesn’t take
place, the agreement-phrase becomes vacuous, and we assume here
that vacuous phrases are not projected at all, but nothing hinges on this
assumption.

When V is plain, there is no position for O and V to move to.

\subsection{Determining PF} 
\label{lasz:sec:23}

Syntax will thus yield a representation where V has copies in the
functional heads it agrees with, and O has copies in the specifiers of
the same functional phrases. In line with \citeapos{Bobaljik.Wurmbrand.2012} 
framework (which builds on \citealp{Bobaljik.1995,Bobaljik.2002,Brody.1995,ErteschikShir.1997}, a.o.), 
this representation is first sent to LF, and
LF decides which copies to interpret, and this will always be the
highest copies. Then possible PF-realizations of this LF are evaluated.
PF-candidates differ in which copies of V to realize (which is trivial,
given the inviolable constraint discussed above) and which copies of
O to realize. We propose the following violable and equally-ranked
constraints for PF to achieve this. Whichever PF-candidate(s) fulfil the most constraints will
be grammatical realizations for this LF.

\subsubsection{Scope Transparency (``ScoT'')}

\citet{Bobaljik.Wurmbrand.2012} investigate the connection between the
scope of quantifiers and word orders in two sets of languages: one set
allows QR and doesn’t allow scrambling (e.g. English), and the other
set allows scrambling and doesn’t allow QR (e.g. German, Japanese).
They argue that these properties are connected in the following way:
all languages allow quantifier movement in the syntax and thus at LF,
but the two sets of languages differ in which copy of the quantifier
they pronounce. Pronouncing the high copy amounts to scrambling
and is not available in English, while it is available in German and
Japanese. To capture the distribution of QR, \citeauthor{Bobaljik.Wurmbrand.2012} introduce the
violable constraint Scope Transparency~\REF{lasz:ex:10}, which, given an LF with
a moved quantifier, favors pronouncing the high copy of the quantifier
(scrambling) over pronouncing the low copy of the quantifier (QR). In
English, where scrambling is not an option for PF (this is an inviolable
constraint), Scope Transparency can be violated, giving rise to QR.

\ea
    \label{lasz:ex:10}
    Scope Transparency (ScoT; \citealp[\pnfmt{373}]{Bobaljik.Wurmbrand.2012}): \\
    If the order of two elements at LF is $A\gg B$, the order at PF is $A\gg B$. \\
    \begin{quote}
    We use the symbol $\gg$ to represent the canonical manifestation
    of hierarchical order at the relevant level: roughly, scope at LF
    [\ldots] and linear precedence at PF [\ldots]. \parentext{ibid.}
    \end{quote}
\z 

\citet{Bobaljik.Wurmbrand.2012} proceed to phenomena where, given
two LF choices and two PF choices, 3 out of the 4 possible
combinations are grammatical. They call this the “3⁄4 signature” and
argue that this is the characteristic pattern of the interaction of ScoT
with another violable constraint. If both constraints favor the same
PF-candidate, this PF will be the “best” PF of the given LF. If the
constraints disagree, however, this results in a tie, and both PF
candidates are the “best” choice. This pattern is illustrated in Table~\ref{lasz:tab:4}.

\begin{table}
    \begin{tabular}{cccccc}
        \lsptoprule 
        LF & PFs & ScoT & Another constraint & \multicolumn{2}{c}{Prediction} \\\midrule
        \laszTabRowA{} & $A \gg B$ & \cmark & \cmark & 2/2 & \HandLeft \\
        \laszTabRowA \multirow{-2}{*}{$A \gg B$} & $B \gg A$ & * & * & 0/2 & \\
        \midrule
        \laszTabRowB & $A \gg B$ & * & \cmark & 1/2 & \HandLeft \\ 
        \laszTabRowB \multirow{-2}{*}{$B \gg A$} & $B \gg A$ & \cmark & * & 1/2 & \HandLeft \\
        \lspbottomrule
    \end{tabular}
    \caption{The 3⁄4 signature \citep[\pnfmt{385}f.]{Bobaljik.Wurmbrand.2012}}
    \label{lasz:tab:4}
\end{table}

Regarding object-shift in ASL and Libras, we assume syntactic
movement of O as explained above, which results in two (or more)
copies of O. Since LF always interprets the highest copy of O, we take
ScoT to be fulfilled when PF realizes the highest copy of O as well.

\subsubsection{Canonical complement order (``CCO'')}

We also adopt \citeauthor{Bobaljik.Wurmbrand.2012}'s second violable constraint called Canonical
complement order (CCO), which captures the “cost” of movement that
is otherwise “free”, i.e., not required for the derivation to succeed, not
feature-driven. \citeauthor{Bobaljik.Wurmbrand.2012} note (p. 390) that the idea for such a constraint is
a fairly old one and used to be called “Case Adjacency” in early
GB-models.
In the present work, this constraint represents favoring the surface
order SVO, which we understand as the underlying order of ASL and
Libras, over the derived order SOV.

\subsubsection{If O values V, O precedes V (``Value'')}

\citet{Napoli.SS.2014} have conducted a survey of the order
of major constituents in 42 sign languages and discovered one
generalization and five tendencies regarding the word order variation
between SVO and SOV. We take our cue from their cross-linguistic
tendency in~\REF{lasz:ex:11} (“Generalization 2”).
\ea 
    \label{lasz:ex:11}
    If an argument affects the phonological shape of V, it precedes V. 
    \citep{Napoli.SS.2014}
\z 
\citeauthor{Napoli.SS.2014} understand “affecting the phonological shape
of V” to include classifier predicates, agreeing verbs, pointing verbs,
spatial verbs, and argument-sensitive verbs. Only plain verbs are
exempt. Our data for ASL and Libras, as presented above, show that
\REF{lasz:ex:11} holds only partially in these languages. We certainly observe that
O interacting with V enables the word order SOV, which aligns with~\REF{lasz:ex:11}. 
However, we also observe a distinction between
classifier-agreement and locus-agreement, which \citeauthor{Napoli.SS.2014} 
bundle together. Let us briefly turn to a discussion
about the differences in syntax between these two types of agreement,
and why -- in ASL -- classifier-agreement is more strict (only OV)
than locus-agreement (OV or VO). With respect to why Libras differs
from ASL here, we remain agnostic at this point, and leave possible
consequences of this seemingly parametric difference to future
research.

We assume that classifier-agreement and locus-agreement both
happen in dedicated phrases as Specifier-Head agreement, such that O
moves to the Specifier and V moves to the Head \citep{Benedicto.Brentari.2004,Kuhn.2016}. 
As illustrated above, we term the phrase
for locus-agreement AgrOP, and the phrase for classifier-agreement
ClassOP. Our proposal is that there is a difference between these two types
of agreement in the direction of feature valuation. For
classifier-agreement, what determines V’s handshape is information
about O’s noun-class. Valuation must thus always proceed from O to
V. For locus-agreement, on the other hand, we argue that valuation can
in principle proceed in either direction. The value of the locus ($a$, $b$, $c$,
\ldots) is not inherent to either O or V; all that matters for agreement is
that they share the same value. In~\REF{lasz:ex:12} we see that the value does not
have to originate from O: V introduces the locus $b$ associated with the
indirect object of \textsc{give}, while the IO itself is a body-anchored sign
that is not assigned a locus (as it might be via IX, eyegaze, head-tilt or
shoulder-shift). We interpret this to mean that the locus-feature on the
IO is unvalued.

\ea 
    \label{lasz:ex:12}
    \textsc{book}, \textsc{father} \textsc{give}\laszDir{a}{b} \textsc{mother}.%
    \footnote{%
        In this example, \textsc{book} is topicalized, as indicated by the comma
        following it.
    }%
    \footnote{%
        Note that the SOV word order is not possible here without locus-marking
        (eyegaze for Libras; shoulder-shift or eyegaze+headtilt for ASL):
        \begin{exe} \ex[*]{\textsc{book}, \textsc{father} \textsc{mother} \textsc{give}\laszDir{a}{b}.} \end{exe}
        We do not go into this further, but we note that there might be two possible reasons
        for this: \textsc{father} and \textsc{mother} might need to be differentiated grammatically to
        ensure that the predicate is non-reversible (which may prohibit the ordering SOV
        according to \citealp{Fischer.1974}), or \textsc{mother} (at least) might need to set the value of the
        locus feature when to the left of V.
    }
    \hfill 
    \cmark ASL,
    \cmark Libras 
    \\
    \parentext{constructed after \citealp[473]{Kuhn.2016}}
\z 

\subsubsection{Phonologically heavier elements come later (``Heavy'')}

Verbs marked with durative aspect have reduplicated movement
\citep{Klima.Bellugi.1979} and are thus phonologically heavier than
non-aspect-marked verbs \citep{Brentari.1998}. However, verbs marked
with lo\-cus-agree\-ment or class\-i\-fier-agree\-ment are not phonologically
heavier \citep{Brentari.1998}.

This is supported by the following data.%
\footnote{%
    In our judgments of ASL, `quickly' is not expressed by a manual sign, so
    we substituted the ASL adverb \LaConcept{SOMETIMES} for the Libras adverb \LaConcept{QUICKLY}.
}
In SVO sentences, adverbs can appear in a position on the left~\REF{lasz:ex:13}, or on the
far right of the clause~\REF{lasz:ex:14}.
\ea 
    \label{lasz:ex:13}
    %  S Adv V O
    \ea 
        \LaConcept{MY} \LaConcept{SISTER} \textbf{\LaConcept{QUICKLY}} \LaConcept{SEND}\laszPlain{} \LaConcept{LETTER}.
        \hfill 
        \cmark Libras
        \\
         ‘My sister quickly sent a/the letter.’
    \ex 
        \LaConcept{MY} \LaConcept{SISTER} \textbf{\LaConcept{SOMETIMES}} \LaConcept{SEND}\laszPlain{} \LaConcept{LETTER}. 
        \hfill 
        \cmark ASL 
        \\ 
        ‘My sister sometimes sends a letter.’
    \z 
\ex 
    \label{lasz:ex:14}
    % S V O Adv
    \ea 
        \LaConcept{MY} \LaConcept{SISTER} \LaConcept{SEND}\laszPlain{} \LaConcept{LETTER} \textbf{\LaConcept{QUICKLY}}.
        \hfill 
        \cmark Libras 
        \\ 
        ‘My sister quickly sent a/the letter.’
    \ex 
        \LaConcept{MY} \LaConcept{SISTER} \LaConcept{SEND}\laszPlain{} \LaConcept{LETTER} \textbf{\LaConcept{SOMETIMES}}.
        \hfill 
        \cmark ASL 
        \\ 
        ‘My sister sometimes sends a letter.’
    \z 
\z 
However, in SOV constructions with aspect-marked V, the adverb
cannot appear on the right.
\ea 
    \label{lasz:ex:15}
    % S Adv O V
    \ea[]{ 
        \LaConcept{MY} \LaConcept{SISTER} \textbf{\LaConcept{QUICKLY}} \LaConcept{LETTER} \LaConcept{SEND}\laszAsp{}.
        \hfill 
        \cmark Libras
        \\ 
        \glt ‘My sister repeatedly sends letters quickly.’}
    \ex[]{ 
        \LaConcept{MY} \LaConcept{SISTER} \textbf{\LaConcept{SOMETIMES}} \LaConcept{LETTER} \LaConcept{SEND}\laszAsp{}.
        \hfill 
        \cmark ASL 
        \\ 
        \glt ‘My sister sometimes sends letters repeatedly.’}
    \z 
\ex 
    \label{lasz:ex:16}
    % *S O V Adv
    \ea[*]{
        \LaConcept{MY} \LaConcept{SISTER} \LaConcept{LETTER} \LaConcept{SEND}\laszAsp{} \textbf{\LaConcept{QUICKLY}}.
        \hfill 
        *Libras
    }
    \ex[*]{
        \LaConcept{MY} \LaConcept{SISTER} \LaConcept{LETTER} \LaConcept{SEND}\laszAsp{} \textbf{\LaConcept{SOMETIMES}}.
        \hfill 
        *ASL
    }
    \z 
\z 
In contrast to aspect-marking, this restriction on the placement of
adverbs does not apply to classifier- \REF{lasz:ex:17} or locus-marking \REF{lasz:ex:18} of the
verb.
\ea 
    \label{lasz:ex:17}
    S O V[hs] Adv
    \ea 
         \LaConcept{IX1 WINE DRINK}\laszHs{\laszClaw} \textbf{\LaConcept{QUICKLY}}.
        \hfill 
        \cmark Libras
    \ex 
        \LaConcept{IX1 WINE DRINK}\laszHs{\laszClaw} \textbf{\LaConcept{SOMETIMES}}. 
        \hfill 
        \cmark ASL 
    \z 
\ex 
    \label{lasz:ex:18}
    S O V[loc] Adv
    \ea 
         \LaConcept{IX1 \LaConcept{HOUSE}}\laszLoc{a} \LaConcept{BUY}\laszLoc{a} \textbf{\LaConcept{QUICKLY}}. 
         \hfill 
         \cmark Libras
    \ex 
         \LaConcept{IX1 \LaConcept{HOUSE}}\laszLoc{a} \LaConcept{BUY}\laszDir{1}{a} \textbf{\LaConcept{SOMETIMES}}.
        \hfill 
        \cmark ASL
    \z 
\z 

\subsection{Our account in action}
\label{lasz:sec:24}

Table~\ref{lasz:tab:5} illustrates the application of these PF constraints for object
shift in ASL and Libras. For any given LF (which differ in the position
of where O and V have moved to), the PF candidates of pronouncing
the sequence as SVO or SOV are evaluated. The columns labelled with a constraint 
show whether this constraint is fulfilled (“\cmark”) or violated (“*”).
The prediction is that all candidates with the highest number of
constraints fulfilled (indicated with the selector symbol “\HandLeft”) will be
acceptable, and all candidates with a lower-than-highest count will be
unacceptable. The final column displays the actual judgment. It 
matches the prediction for all rows. 

Table~\ref{lasz:tab:5} illustrates the different outcomes when a plain verb
with no classifier-agreement is used, compared to verbs with
classifier-agreement. The first two rows present the LF without
agreement, i.e. the verb has plain inflection. Since there are no other
copies of O and V, only one PF-candidate is available, namely the one
with the ordering SVO.

The second two rows present the minimally different LF where
classifier-agree\-ment did happen. We see two PF-candidates: the first
one spells out the copy of O inside the VP, and the second one spells
out the copy of O in Spec,Cl\textsubscript{O}P. Note that we do not need to consider
PF-candidates that spell out the higher copy of V, since there is an
inviolable constraint requiring lower V to be spelled out. The
SVO-candidate does not fulfill ScoT, but fulfills CCO. The
SOV-candidate fulfills ScoT, but does not fulfill CCO. We see the 3⁄4
signature arising from the interaction of two violable constraints.

\begin{table}
    \begin{tabular}{ll cccc}
        \lsptoprule 
        LF & PF candidates & ScoT & CCO & Pred. & Judg. \\\midrule
        \multicolumn{2}{l}{
            S \laszLB{VP} V\textsubscript{\laszPlain} O ] } \\
        &   S \laszLB{VP} V\textsubscript{\laszPlain} O ] &
            \cmark & \cmark & \HandLeft & ok \\
        &   S \laszLB{?} O \laszLB{VP} V\textsubscript{\laszPlain} \lasztO{} ] ] &
            \multicolumn{3}{l}{Not applicable\footnotemark{}} & * \\ 
        \midrule
        \multicolumn{2}{l}{
            S \laszLB{ClassOP} O V\textsubscript{\laszHs{\_}} \laszLB{VP} \lasztV{} \lasztO{} ] ] } \\
        &   S \laszLB{ClassOP} \lasztO{} \lasztV{} \laszLB{VP} V\textsubscript{\laszHs{\_}} O  ] ] & 
            * & \cmark & \HandLeft & ok \\ 
        &   S \laszLB{ClassOP} O \lasztV{} \laszLB{VP} V\textsubscript{\laszHs{\_}} \lasztO{} ] ] & 
            \cmark & * & \HandLeft & ok \\ 
        \lspbottomrule
    \end{tabular}
    \caption{No agreement vs. classifier-agreement in Libras}
    \label{lasz:tab:5}
\end{table}
\footnotetext{%
    There is no higher copy of O at LF that PF could spell out. This
    PF-candidate does not exist and is only given for completeness to illustrate that the
    ordering SOV\laszPlain{} is ungrammatical.
}

Table~\ref{lasz:tab:6} illustrates classifier-agreement in ASL. Here
the constraint “Value” comes into play. It
differentiates between the two possible PFs of the second LF, allowing
only SOV ordering in ASL with classifier-agreeing V.

\begin{table}
    \begin{tabular}{ll ccccc}
        \lsptoprule 
        LF & PF candidates & ScoT & CCO & Value & Pred. & Judg. \\\midrule
        \multicolumn{2}{l}{
            S \laszLB{VP} V\textsubscript{\laszPlain} O ] } \\
        &   S \laszLB{VP} V\textsubscript{\laszPlain} O ] &
            \cmark & \cmark & \cmark & \HandLeft 
            & ok \\
        &   S \laszLB{?} O \laszLB{VP} V\textsubscript{\laszPlain} \lasztO{} ] ] &
            \multicolumn{4}{l}{Not applicable} 
            & * \\ 
        \midrule 
        \multicolumn{2}{l}{
            S \laszLB{ClassOP} O V\textsubscript{\laszHs{\_}} \laszLB{VP} \lasztV{} \lasztO{} ] ] } \\ 
        &   S \laszLB{ClassOP} \lasztO{} \lasztV{} \laszLB{VP} V\textsubscript{\laszHs{\_}} O ] ] & 
            * & \cmark & * &  
            & * \\ 
        &   S \laszLB{ClassOP} O \lasztV{} \laszLB{VP} V\textsubscript{\laszHs{\_}} \lasztO{} ] ] & 
            \cmark & * & \cmark & \HandLeft 
            & ok \\ 
        \lspbottomrule 
    \end{tabular}
    \caption{No agreement vs. classifier-agreement in ASL}
    \label{lasz:tab:6}
\end{table}

Locus-agreement in both ASL and Libras mirrors Table~\ref{lasz:tab:5}
(because, as argued above, ``Value'' does not apply to
locus-agreement), except that it is less optional for O. I.e., whenever O
is associated with a locus and V allows locus-agreement with O, V
will agree (\citealp{Padden.1988}; but see the corpus studies by 
\citealp{DeBeuzeville.etal.2009} for Auslan, 
\citealp{Costello.2015} for LSE, and 
\citealp{Fenlon.etal.2018} for BSL 
showing that while locus-agreement on V with O is more
frequent than with S, this is by no means obligatory in these three sign
languages).

Aspect-marking on V is also not optional, given that it is the
only marking of aspectual information in a sentence. V that is marked
with durational aspect modifies its movement to contain a circular
motion, which corresponds to phonological heaviness \citep{Brentari.1998}.
Thus, the constraint ``Heavy'', which was
fulfilled for all PF-candidates in Tables~\ref{lasz:tab:5} and~\ref{lasz:tab:6}, differentiates the
PF-candidates in Table~\ref{lasz:tab:7}:

\begin{table}
    \begin{tabular}{ll ccccc}
        \lsptoprule
        LF & PF candidates & ScoT & CCO & Heavy & Pred. & Judg. \\ 
        \midrule
        \multicolumn{2}{l}{
            S \laszLB{AspP} O V\textsubscript{\laszAsp{}} \laszLB{VP} \lasztV{} \lasztO{} ] ] } \\ 
        &   S \laszLB{AspP} \lasztO{} \lasztV{} \laszLB{VP} V\textsubscript{\laszAsp} O ] ] &
            * & \cmark & * & &
            *\footnotemark{} \\ 
        &   S \laszLB{AspP} O \lasztV{} \laszLB{VP} V\textsubscript{\laszAsp} \lasztO{} ] ] &
            \cmark & * & \cmark & \HandLeft &
            ok \\ 
        \lspbottomrule
    \end{tabular}
    \caption{Aspect-marking in ASL and Libras}
    \label{lasz:tab:7}
\end{table}
\footnotetext{
    \citet{Braze.2004} reports that his consultants like the order SV\laszAsp{}O and
    dislike the order SOV\laszAsp{}. This may be due to an influence of English on ASL, as
    \citet{Matsuoka.1997} notes, or it could be dialectal variation within ASL, as \citeauthor{Braze.2004}
    suspects, or both.
}

% ===================================================================
\section{Problems with previous proposals}
\label{lasz:sec:3}

It has been proposed that the obligatory word order change from VO
to OV when V carries aspect-marking is due to V moving into a
right-headed Aspect phrase to collect this extra morphology. \citet{Fischer.Janis.1992}
suggested such an analysis, although for verb-doubling
constructions of the form SV\laszPlain{}OV\laszAsp{}. This proposal has then
been picked up, refined, and applied to object shift constructions by
\citet{Matsuoka.1997} and \citet{Braze.2004}.
A verb-doubling construction is illustrated in~\REF{lasz:ex:19}.
\ea 
    \label{lasz:ex:19}
     \LaConcept{SALLY} \LaConcept{TYPE}\laszPlain{} \LaConcept{HER} \LaConcept{TERM} \LaConcept{PAPER} \LaConcept{TYPE}\laszAsp{}.
     \hfill 
     \cmark ASL, \cmark Libras 
     \\ 
     \glt ‘Sally typed her term paper for a very long time.’
     \citep[from][]{Fischer.Janis.1992}
\z 
The idea is that the two verbs, which always share the same root, are
two realizations of the same chain. The first verb appears in its base
position inside VP. The second verb carries additional marking, thus
appears in a position with additional morphology, which is the Asp
head, which must be on the right given this construction’s word order.
\citet{Matsuoka.1997} proposes the derivation in Fig.~\ref{lasz:fig:2}. In a verb-doubling
construction, both copies of V are pronounced: the copy inside VP in
its plain form, and the copy in Asp with aspect-marking.

\begin{figure}
    \begin{tikzpicture}
        \Tree 
        [.IP
            [.AspP
                [.VP
                    [.\phantom{X} ]
                    [.V'
                        \node (Vbase) {V};
                        [.O ]
                    ]
                ]
                [.Asp
                    \node (Vmoved) {V};
                    [.affix ]
                ]
            ]
        ]
        \draw[->] (Vbase) ..controls +(south east:1.5) and +(south:1) .. (Vmoved);
    \end{tikzpicture}
    \caption{\citeapos{Matsuoka.1997} analysis of aspect-marked verbs}
    \label{lasz:fig:2}
\end{figure}

We reject this analysis on two grounds: one empirical and one
theory-internal, which we discuss in turn in subsections~\ref{lasz:sec:32} and~\ref{lasz:sec:33}.


\subsection{Adverb placement}
\label{lasz:sec:32}

The empirical problem with the right-headed-Asp analysis is that it
cannot capture the following data for adverb placement, while our PF-based account can.

In general, adverbs can appear to the left of V (temporal
adverbs before S, as illustrated in~\REF{lasz:ex:20}; adverbs of manner after S, as
illustrated in~\REF{lasz:ex:21}) or to the right of V (in sentence-final position,
\REF{lasz:ex:22}). We remain agnostic as to whether this is derived via optional
left- or right-adjunction of the adverb to TP resp. \textit{v}P, or via
extraposition of the adverb.
\ea 
    \label{lasz:ex:20}
    Temporal adverbs adjoin on the left.
    \ea 
        \textbf{\LaConcept{YESTERDAY}} \LaConcept{MY} \LaConcept{SISTER} \LaConcept{SEND}\laszPlain{} \LaConcept{LETTER}.
        \hfill 
        \cmark Libras, \cmark ASL 
    \ex 
        \textbf{\LaConcept{YESTERDAY}} \LaConcept{MY} \LaConcept{SISTER} \LaConcept{LETTER} \LaConcept{SEND}\laszAsp{}.
        \hfill 
        \cmark Libras, \cmark ASL 
        \\ 
        both: ‘Yesterday my sister sent a/the letter(s).’
    \z 
\ex 
    \label{lasz:ex:21}
    Frequency adverbs adjoin on the left.
    \ea 
        \LaConcept{MY} \LaConcept{SISTER} \textbf{\LaConcept{QUICKLY}} \LaConcept{SEND}\laszPlain{} \LaConcept{LETTER}.
        \hfill 
        \cmark Libras 
    \ex 
        \LaConcept{MY} \LaConcept{SISTER} \textbf{\LaConcept{QUICKLY}} \LaConcept{LETTER} \LaConcept{SEND}\laszAsp{}.
        \hfill 
        \cmark Libras 
        \\ 
        both: ‘My sister quickly sent a/the letter(s).’
    \ex 
        \LaConcept{MY} \LaConcept{SISTER} \textbf{\LaConcept{SOMETIMES}} \LaConcept{LETTER} \LaConcept{SEND}\laszDir{}{1}.
        \hfill 
        \cmark ASL 
        \\ 
        ‘Sometimes my sister sends me a letter.’
    \z 
\ex 
    \label{lasz:ex:22}
    Adverbs adjoin on the right in SVO clauses.
    \ea 
        \LaConcept{MY} \LaConcept{SISTER} \LaConcept{SEND}\laszPlain{} \LaConcept{LETTER} \textbf{\LaConcept{YESTERDAY}}.
        \hfill 
        \cmark Libras, \cmark ASL 
    \ex 
        \LaConcept{MY} \LaConcept{SISTER} \LaConcept{SEND}\laszPlain{} \LaConcept{LETTER} \textbf{\LaConcept{QUICKLY}}.
        \hfill 
        \cmark Libras 
    \ex 
        \LaConcept{MY} \LaConcept{SISTER} \LaConcept{LETTER} \LaConcept{SEND}\laszDir{}{1} \textbf{\LaConcept{SOMETIMES}}.
        \hfill 
        \cmark ASL 
    \z 
\z 
However, in object-shift clauses with aspect-marked V, adverbs must
adjoin on the left. They do not have the option of appearing in
sentence-final position. This is illustrated in~\REF{lasz:ex:23}.
\ea 
    \label{lasz:ex:23}
    Adverbs cannot adjoin on the right in SOV\laszAsp{} clauses.
    \ea[*]{
        \LaConcept{MY} \LaConcept{SISTER} \LaConcept{LETTER} \LaConcept{SEND}\laszAsp{} \textbf{\LaConcept{YESTERDAY}}.
        \hfill
        *Libras, *ASL
    }
    \ex[*]{
        \LaConcept{MY} \LaConcept{SISTER} \LaConcept{LETTER} \LaConcept{SEND}\laszAsp{} \textbf{\LaConcept{QUICKLY}}.
        \hfill 
        *Libras
    }
    \z 
\z 

That the examples in~\REF{lasz:ex:23} should be ungrammatical is completely
unexpected under the theories proposed by \citet{Fischer.Janis.1992},
\citet{Matsuoka.1997}, and \citet{Braze.2004}. Why should movement of V to
Asp prohibit an adverb adjoining on the right, or extraposing? In
contrast to this, our account captures these data perfectly. Since we
postulate a PF-constraint stating that phonologically heavy V should
be in sentence-final position, SOV-constructions with V\laszAsp{} followed
by another element are penalized.%
\footnote{%
    Elements that are typically in sentence-final position and may appear
    even after V\laszAsp{} include the subject pronoun copy. This is outside the scope of this
    paper, but it seems to us that copying a subject pronoun into the sentence-final
    position might be an inviolable constraint, taking precedence over the constraints
    discussed here.
}
Our derivations are illustrated in Tables~\ref{lasz:tab:9} and~\ref{lasz:tab:10}.

\begin{table}
    \begin{tabular}{ll cccc}
        \lsptoprule
        LF & PF candidates & ScoT & CCO & Pred. & Judg. \\\midrule
        \multirow{4}{*}{
            S \laszLB{VP} V O ]}
        & S \laszLB{} Adv \laszLB{VP} V\textsubscript{\laszPlain} O 
            & \cmark & \cmark & \HandLeft & ok \\
        & S \laszLB{} \laszLB{VP} V\textsubscript{\laszPlain} O ] Adv ] 
            & \cmark & \cmark & \HandLeft & ok \\ 
        & S \laszLB{} Adv \laszLB{?} O \laszLB{VP} V\textsubscript{\laszPlain} 
            & \multicolumn{3}{c}{Not applicable} & * \\ 
        & S \laszLB{} \laszLB{?} O \laszLB{VP} V\textsubscript{\laszPlain} ] ] Adv ] 
            & \multicolumn{3}{c}{Not applicable} & * \\ 
        \lspbottomrule 
    \end{tabular}
    \caption{Adverb-placement options with plain V}
    \label{lasz:tab:9}
\end{table}

\begin{table}
    \begin{tabular}{ll ccccc}
        \lsptoprule 
        LF & PF candidates & ScoT & CCO & Heavy & Pred. & Judg. \\\midrule 
        \multicolumn{2}{l}{
            S \laszLB{AspP} O V \laszLB{VP} \lasztV{} \lasztO{} ] ] } \\
        & S [ Adv \laszLB{AspP} \lasztO{} \lasztV{} \laszLB{VP} V O ] ] ] & 
            * & \cmark & * & & * \\ 
        & S [ \laszLB{AspP} \lasztO{} \lasztV{} \laszLB{VP} V O ] ] Adv ] & 
            * & \cmark & * & & * \\ 
        & S [ Adv \laszLB{AspP} O \lasztV{} \laszLB{VP} V \lasztO{} ] ] ] & 
            \cmark & * & \cmark & \HandLeft & ok \\ 
        & S [ \laszLB{AspP} O \lasztV{} \laszLB{VP} V \lasztO{} ] ] Adv ] & 
            \cmark & * & * & & * \\
        \lspbottomrule
    \end{tabular}
    \caption{Adverb-placement options with aspect-marked V}
    \label{lasz:tab:10}
\end{table}

Our account also correctly predicts that when V is marked with
classifier- or locus-agreement, adverb placement is possible on the left
as well as on the right. Consider the data in \REF{lasz:ex:24}--\REF{lasz:ex:27} and our
derivations in Table~\ref{lasz:tab:11} (given for locus-agreement only;
classifier-agreement is completely parallel in Libras, and in ASL
the constraint ``Value'' comes into play additionally).

\ea 
    \label{lasz:ex:24}
    Adverb placement with classifier-agreement in Libras
    \ea 
        \LaConcept{IX1} \textbf{\LaConcept{QUICKLY}} \LaConcept{DRINK}\laszHs{\laszClaw} \LaConcept{WINE}. 
        
    \ex 
        \LaConcept{IX1 DRINK}\laszHs{\laszClaw} \LaConcept{WINE} \textbf{\LaConcept{QUICKLY}}.
        \hfill \cmark Libras
    \ex 
        \LaConcept{IX1} \textbf{\LaConcept{QUICKLY}} \LaConcept{WINE DRINK}\laszHs{\laszClaw}.
        \hfill \cmark Libras
    \ex 
        \LaConcept{IX1 WINE DRINK}\laszHs{\laszClaw} \textbf{\LaConcept{QUICKLY}}. 
        \hfill \cmark Libras
        \\ 
        \glt all: ‘I quickly drank wine.’
    \z 
\ex 
    \label{lasz:ex:25}
    Adverb placement with classifier-agreement in ASL
    \ea 
        \LaConcept{IX1} \textbf{\LaConcept{SOMETIMES}} \LaConcept{DRINK}\laszHs{\laszClaw} \LaConcept{WINE}.
        \hfill \cmark ASL
    \ex 
        \LaConcept{IX1 DRINK}\laszHs{\laszClaw} \LaConcept{WINE} \textbf{\LaConcept{SOMETIMES}}.
        \hfill \cmark ASL
    \ex 
        \LaConcept{IX1} \textbf{\LaConcept{SOMETIMES}} \LaConcept{WINE DRINK}\laszHs{\laszClaw}.
        \hfill \cmark ASL
    \ex 
        \LaConcept{IX1 WINE DRINK}\laszHs{\laszClaw} \textbf{\LaConcept{SOMETIMES}}. 
        \hfill \cmark ASL
        \\ 
        \glt all: ‘I drink wine sometimes.’
    \z 
\ex 
    \label{lasz:ex:26}
    Adverb placement with locus-agreement in Libras
    \ea 
        \LaConcept{IX1} \textbf{\LaConcept{QUICKLY}} \LaConcept{BUY}\laszLoc{a} \LaConcept{HOUSE}\laszLoc{a}.
        \hfill \cmark Libras
    \ex 
        \LaConcept{IX1} \LaConcept{BUY}\laszLoc{a} \LaConcept{HOUSE}\laszLoc{a} \textbf{\LaConcept{QUICKLY}}.
        \hfill \cmark Libras
    \ex 
        \LaConcept{IX1} \textbf{\LaConcept{QUICKLY}} \LaConcept{HOUSE}\laszLoc{a} \LaConcept{BUY}\laszLoc{a}.
        \hfill \cmark Libras
    \ex 
        \LaConcept{IX1 \LaConcept{HOUSE}}\laszLoc{a} \LaConcept{BUY}\laszLoc{a} \textbf{\LaConcept{QUICKLY}}. 
        \hfill \cmark Libras
        \\ 
        \glt all: ‘I quickly bought the house.’
    \z 
\ex
    \label{lasz:ex:27}
    Adverb placement with locus-agreement in ASL
    \ea 
        \LaConcept{IX1} \textbf{\LaConcept{SOMETIMES}} \LaConcept{BUY}\laszDir{1}{a} \LaConcept{HOUSE}\laszLoc{a}.
        \hfill \cmark ASL
    \ex 
        \LaConcept{IX1} \LaConcept{BUY}\laszDir{1}{a} \LaConcept{HOUSE}\laszLoc{a} \textbf{\LaConcept{SOMETIMES}}.
        \hfill \cmark ASL
    \ex 
        \LaConcept{IX1} \textbf{\LaConcept{SOMETIMES}} \LaConcept{HOUSE}\laszLoc{a} \LaConcept{BUY}\laszDir{1}{a}.
        \hfill \cmark ASL
    \ex 
        \LaConcept{IX1 \LaConcept{HOUSE}}\laszLoc{a} \LaConcept{BUY}\laszDir{1}{a} \textbf{\LaConcept{SOMETIMES}}. 
        \hfill \cmark ASL
        \\ 
        all: `I sometimes buy a house.' 
    \z 
\z 

\begin{table}
    \begin{tabular}{ll cccc}
        \lsptoprule 
        LF & PF candidates & ScoT & CCO & Pred. & Judg. \\\midrule
        \multicolumn{2}{l}{
            S \laszLB{AgrOP} O V \laszLB{VP} \lasztV{} \lasztO{} ] ] } \\ 
        & S [ Adv \laszLB{AgrOP} \lasztO{} \lasztV{} \laszLB{VP} V O ] ] ] &
            * & \cmark & \HandLeft & ok \\ 
        & S [ \laszLB{AgrOP} \lasztO{} \lasztV{} \laszLB{VP} V O ] ] Adv ] &
            * & \cmark & \HandLeft & ok \\ 
        & S [ Adv \laszLB{AgrOP} O \lasztV{} \laszLB{VP} V \lasztO{} ] ] ] &
            \cmark & * & \HandLeft & ok \\ 
        & S [ \laszLB{AgrOP} O \lasztV{} \laszLB{VP} V \lasztO{} ] ] Adv ] &
            \cmark & * & \HandLeft & ok \\ 
        \lspbottomrule
    \end{tabular}
    \caption{Adverb-placement options with locus-marked V}
    \label{lasz:tab:11}
\end{table}

\subsection{Mixed headedness}
\label{lasz:sec:33}

The conceptual problem with the proposal that V moves from a
left-headed VP to a right-headed AspP is that it violates the
Final-over-Final Constraint (FOFC, pronounced [fofk]; see 
\citealp{BiberauerHR.2014,SheehanBRH.2017}, and much subsequent work). The
FOFC is an implicational universal over the headedness of phrases
within an extended projection. There are 2 types of extended
projections: the verbal one containing V, \textit{v}, Asp, Aux, T, C, etc; and
the nominal one containing N, Num, Poss, D, etc. The implicational
universal FOFC states informally:
\ea 
    \label{lasz:ex:28}
     If a phrase in an extended projection is right-headed, all
    phrases below this phrase in this extended projection are also
    right-headed. \citep[after][]{BiberauerHR.2014}
\z 
The logical consequence of~\REF{lasz:ex:28} is~\REF{lasz:ex:29}:
\ea 
    \label{lasz:ex:29}
    If a phrase in an extended projection is left-headed, all
    phrases above this phrase in this extended projection are also
    left-headed.
\z 
Thus, if in a language, CP is right-headed, then TP, \textit{v}P, VP, and all
other verbal projections below CP are also right-headed (Fig.~\ref{lasz:fig:3a} vs.~\ref{lasz:fig:3d}
for $Y$=C). If VP is right-headed, this does not make a statement about
the headedness of the next higher verbal projection (Fig.~\ref{lasz:fig:3a} vs.~\ref{lasz:fig:3c} for
$X$=V). If VP is left-headed, then \textit{v}P, TP, CP, or any other verbal
projection higher than VP cannot be right-headed because this would
violate~\REF{lasz:ex:29} (Fig.~\ref{lasz:fig:3b} vs.~\ref{lasz:fig:3d} for $X$=V). If CP is left-headed, any
headedness is possible for the next lower verbal projection (Fig.~\ref{lasz:fig:3b}
and~\ref{lasz:fig:3c} for $Y$=C). It is easy to see that the proposal that in ASL, VP is
left-headed and AspP is right-headed, is incompatible with the FOFC
(Fig.~\ref{lasz:fig:3d} for $X$=V and $Y$=Asp).

\begin{figure}
    \centering
    \begin{subfigure}{.2\textwidth}
        \centering 
        \begin{tikzpicture}
            \Tree 
            [.\node {\cmark $Y$P\phantom{\cmark}};
                [.\node {$X$P};
                    \node {\ldots};
                    \node {$X$};
                ]
                \node {$Y$};
            ]
        \end{tikzpicture}
        \caption{Final over final}
        \label{lasz:fig:3a}
    \end{subfigure}
    \hfill 
    \begin{subfigure}{.25\textwidth}
        \centering 
        \begin{tikzpicture}
            \Tree 
            [.\node {\cmark $Y$P\phantom{\cmark}};
                \node {$Y$};
                [.\node {$X$P};
                    \node {$X$};
                    \node {\ldots};
                ]
            ]
        \end{tikzpicture}
        \caption{Initial over initial}
        \label{lasz:fig:3b}
    \end{subfigure}
    \hfill 
    \begin{subfigure}{.2\textwidth}
        \centering 
        \begin{tikzpicture}
            \Tree 
            [.\node {\cmark $Y$P\phantom{\cmark}};
                \node {$Y$};
                [.\node {$X$P};
                    \node {\ldots};
                    \node {$X$};
                ]
            ]
        \end{tikzpicture}
        \caption{Initial over final}
        \label{lasz:fig:3c}
    \end{subfigure}
    \hfill 
    \begin{subfigure}{.25\textwidth}
        \centering 
        \begin{tikzpicture}
            \Tree 
            [.\node {*$Y$P\phantom{*}};
                [.\node {$X$P};
                    \node {$X$};
                    \node {\ldots};
                ]
                \node {$Y$};
            ]
        \end{tikzpicture}
        \caption{*Final over initial\footnotemark{}}
        \label{lasz:fig:3d}
    \end{subfigure}
    \caption{Possible and impossible headedness combinations}
    \label{lasz:fig:3}
\end{figure}
\footnotetext{%
    The FOFC is unintuitively named, as the configuration Final-over-final in Fig.~\ref{lasz:fig:3a} does not stand out in the paradigm. 
    However, the name can be understood as: \laszLB{Topic} Final ] over \laszLB{Focus} Final ] Constraint, 
    saying that if $Y$P is head-final, $X$P must be head-final as well. 
}

The FOFC is well-motivated based on a variety of synchronic
and diachronic cross-linguistic studies, connected to other linguistic
universals, finds support in processing and language acquisition, and
fits in well with current assumptions in Minimalism (see arguments
provided in \citealp{BiberauerHR.2014,SheehanBRH.2017}, and many other
works). We believe that an account that is compatible with the FOFC
should on conceptual grounds be preferred over an account that
violates it.

% ===================================================================
\section{Conclusions}
\label{lasz:sec:4}

In this paper, we have proposed a novel account for object shift in
ASL and Libras. Like accounts of object shift in spoken languages
(\citealp{Holmberg:1986} and many others), and expanding the proposals for
ASL \& Libras by \citet{Quadros.1999,Quadros.etal.2004,Quadros.DLM.2010,Gokgoz.2013}, 
we assume that the reversal
of V and O happens in the syntax by overt movement of O. We have
amended this syntactic proposal by violable constraints at PF,
following \citeapos{Bobaljik.Wurmbrand.2012} framework, which
determine the copy of O that will be realized at PF. This captures the
object-shift paradigm, whereby object-shift is always triggered by
aspect-marking, always triggered by classifier-agreement in ASL,
optionally triggered by classifier-agreement in Libras, optionally
triggered by locus-agreement in both languages, and unacceptable
otherwise. Furthermore, we have shown that our theory makes better
predictions than the competing proposal by 
\citet{Fischer.Janis.1992,Matsuoka.1997,Braze.2004}, 
who argue that the reversal of V
and O happens by movement of V into a right-headed Aspect phrase.
We have shown data from adverb placement that is problematic for
their proposal but derived perfectly in our theory; and we have given a
theory-internal argument against assuming a right-headed Aspect
phrase dominating a left-headed VP -- namely that it violates the
implicational universal FOFC.

In this way, we have provided additional evidence for the
number of verbal projections in ASL and Libras, thereby for the size
of the VP-domain, and for the importance of considering
PF-constraints in the investigation of syntactic phenomena.

% ===================================================================
\section*{Acknowledgements}

We are grateful to Rosie Noschese for providing additional judgments
for ASL. We thank Kadir Gökgöz and two anonymous reviewers for
helpful discussion, and Debbie Chen Pichler, Martina Laszakovits, and
Nora Ungar for their help in acquiring literature. 

% ===================================================================
\printbibliography[heading=subbibliography,notkeyword=this]

\end{document}
